\section{Conclusion}
\label{sec:conclusion}

We began by observing a puzzling gap: hallucination detection methods report AUROC 0.70--0.85 in isolation, yet practitioners lack guidance on which method works for their specific benchmark. Our pilot study reveals why---method performance is benchmark-sensitive in ways previous evaluations could not expose.

\textbf{Summary:}
\begin{enumerate}
    \item Benchmark sensitivity is real (SE: 0.551 HaluEval vs 0.289 TruthfulQA)
    \item Self-consistency requires different similarity metrics
    \item Controlled comparison is essential for method selection
\end{enumerate}

\textbf{Future Directions:}
\begin{itemize}
    \item Investigate TruthfulQA inversion (labeling methodology vs model-specific effects)
    \item Replace BERTScore with NLI-based consistency for self-consistency
    \item Full-scale evaluation (N=817+ samples) with multiple seeds
\end{itemize}

Our findings suggest that hallucination detection is harder than published results imply---not because methods don't work, but because benchmark choice fundamentally affects evaluation outcomes. We hope this work encourages more rigorous, multi-benchmark evaluation practices.
